
\subsubsection{PID Control}


In 1922, Nicolas Minorsky presented a very clear analysis of the control actions necessary to provide effective control of a system whose exact dynamics were un- known. He analyzed the actions taken by a good helmsman steering a ship and translated these actions into the appropriate mathematical formulations \parencite{bennett1993pid}. It is now know as the Proportional-Integral-Derivative (PID) controller. PID is a type of controller that is most commonly used and applied in mechanical systems \parencite{idrissi_review_2022}.

The controller consists of three main components:

\begin{itemize}
    \item Proportional ($K_p$): Acts like a virtual spring that pulls the system toward the desired position; the control action increases as the error gets larger.
    \item Integral ($K_i$): Accumulates the error over time to eliminate steady-state error, ensuring the drone reaches the exact target even in the presence of disturbances like gravity.
    \item erivative ($K_d$): Acts as a virtual damper to reduce oscillations and prevent the system from overshooting its target.
\end{itemize}

The idealized control law is expressed as:

$$u(t) = K_p e(t) + K_i \int e(t) dt + K_d \frac{de(t)}{dt}$$
